\documentclass[11pt]{article}
\newcommand{\ListaTitulo}{Gabarito --- Lista 12}
% Preâmbulo compartilhado — MATD48, Listas de Exercícios 2026
% Usado por Lista{NN}.tex e Gabarito{NN}.tex via % Preâmbulo compartilhado — MATD48, Listas de Exercícios 2026
% Usado por Lista{NN}.tex e Gabarito{NN}.tex via % Preâmbulo compartilhado — MATD48, Listas de Exercícios 2026
% Usado por Lista{NN}.tex e Gabarito{NN}.tex via \input{preamble}

\usepackage[utf8]{inputenc}
\usepackage[T1]{fontenc}
\usepackage[brazil]{babel}
\usepackage{fullpage}
\usepackage{amsmath,amssymb}
\usepackage{enumitem}
\usepackage{xcolor}
\usepackage{fancyhdr}
\usepackage{titlesec}
\usepackage{listings}
\usepackage{hyperref}
\usepackage{booktabs}
\usepackage{graphicx}

\definecolor{matd48azul}{HTML}{0B4F6C}
\definecolor{matd48verde}{HTML}{2E8B57}

\titleformat{\section}{\normalfont\Large\bfseries\color{matd48azul}}{\thesection}{1em}{}
\titleformat{\subsection}{\normalfont\large\bfseries\color{matd48azul}}{\thesubsection}{1em}{}

\providecommand{\ListaTitulo}{Lista de Exercícios}

\setlength{\headheight}{14pt}
\pagestyle{fancy}
\fancyhf{}
\lhead{\small MATD48 --- Planejamento de Experimentos A}
\rhead{\small \ListaTitulo}
\cfoot{\thepage}
\renewcommand{\headrulewidth}{0.4pt}

\lstset{
  basicstyle=\ttfamily\small,
  breaklines=true,
  frame=single,
  columns=fullflexible,
  backgroundcolor=\color{gray!8},
  commentstyle=\color{matd48verde},
  keywordstyle=\color{matd48azul}\bfseries,
  language=R,
  extendedchars=true,
  literate=%
    {á}{{\'a}}1 {à}{{\`a}}1 {â}{{\^a}}1 {ã}{{\~a}}1
    {é}{{\'e}}1 {ê}{{\^e}}1 {í}{{\'i}}1 {ó}{{\'o}}1
    {ô}{{\^o}}1 {õ}{{\~o}}1 {ú}{{\'u}}1 {ç}{{\c c}}1
    {Á}{{\'A}}1 {À}{{\`A}}1 {Â}{{\^A}}1 {Ã}{{\~A}}1
    {É}{{\'E}}1 {Ê}{{\^E}}1 {Í}{{\'I}}1 {Ó}{{\'O}}1
    {Ô}{{\^O}}1 {Õ}{{\~O}}1 {Ú}{{\'U}}1 {Ç}{{\c C}}1
}

\hypersetup{colorlinks=true, linkcolor=matd48azul, urlcolor=matd48azul, citecolor=matd48azul}

\newcommand{\cabecalho}[2]{%
  \begin{center}
    {\Large\bfseries\color{matd48azul} #1}\\[0.3em]
    {\large #2}\\[0.3em]
    \rule{\linewidth}{0.6pt}
  \end{center}
  \vspace{0.5em}
}

\newenvironment{questao}[1]{\vspace{1em}\noindent\textbf{Questão #1.}\hspace{0.4em}}{\par}
\newenvironment{solucao}{\vspace{0.3em}\noindent\textit{Solução.}\hspace{0.4em}\color{black!85}}{\par\vspace{0.5em}}


\usepackage[utf8]{inputenc}
\usepackage[T1]{fontenc}
\usepackage[brazil]{babel}
\usepackage{fullpage}
\usepackage{amsmath,amssymb}
\usepackage{enumitem}
\usepackage{xcolor}
\usepackage{fancyhdr}
\usepackage{titlesec}
\usepackage{listings}
\usepackage{hyperref}
\usepackage{booktabs}
\usepackage{graphicx}

\definecolor{matd48azul}{HTML}{0B4F6C}
\definecolor{matd48verde}{HTML}{2E8B57}

\titleformat{\section}{\normalfont\Large\bfseries\color{matd48azul}}{\thesection}{1em}{}
\titleformat{\subsection}{\normalfont\large\bfseries\color{matd48azul}}{\thesubsection}{1em}{}

\providecommand{\ListaTitulo}{Lista de Exercícios}

\setlength{\headheight}{14pt}
\pagestyle{fancy}
\fancyhf{}
\lhead{\small MATD48 --- Planejamento de Experimentos A}
\rhead{\small \ListaTitulo}
\cfoot{\thepage}
\renewcommand{\headrulewidth}{0.4pt}

\lstset{
  basicstyle=\ttfamily\small,
  breaklines=true,
  frame=single,
  columns=fullflexible,
  backgroundcolor=\color{gray!8},
  commentstyle=\color{matd48verde},
  keywordstyle=\color{matd48azul}\bfseries,
  language=R,
  extendedchars=true,
  literate=%
    {á}{{\'a}}1 {à}{{\`a}}1 {â}{{\^a}}1 {ã}{{\~a}}1
    {é}{{\'e}}1 {ê}{{\^e}}1 {í}{{\'i}}1 {ó}{{\'o}}1
    {ô}{{\^o}}1 {õ}{{\~o}}1 {ú}{{\'u}}1 {ç}{{\c c}}1
    {Á}{{\'A}}1 {À}{{\`A}}1 {Â}{{\^A}}1 {Ã}{{\~A}}1
    {É}{{\'E}}1 {Ê}{{\^E}}1 {Í}{{\'I}}1 {Ó}{{\'O}}1
    {Ô}{{\^O}}1 {Õ}{{\~O}}1 {Ú}{{\'U}}1 {Ç}{{\c C}}1
}

\hypersetup{colorlinks=true, linkcolor=matd48azul, urlcolor=matd48azul, citecolor=matd48azul}

\newcommand{\cabecalho}[2]{%
  \begin{center}
    {\Large\bfseries\color{matd48azul} #1}\\[0.3em]
    {\large #2}\\[0.3em]
    \rule{\linewidth}{0.6pt}
  \end{center}
  \vspace{0.5em}
}

\newenvironment{questao}[1]{\vspace{1em}\noindent\textbf{Questão #1.}\hspace{0.4em}}{\par}
\newenvironment{solucao}{\vspace{0.3em}\noindent\textit{Solução.}\hspace{0.4em}\color{black!85}}{\par\vspace{0.5em}}


\usepackage[utf8]{inputenc}
\usepackage[T1]{fontenc}
\usepackage[brazil]{babel}
\usepackage{fullpage}
\usepackage{amsmath,amssymb}
\usepackage{enumitem}
\usepackage{xcolor}
\usepackage{fancyhdr}
\usepackage{titlesec}
\usepackage{listings}
\usepackage{hyperref}
\usepackage{booktabs}
\usepackage{graphicx}

\definecolor{matd48azul}{HTML}{0B4F6C}
\definecolor{matd48verde}{HTML}{2E8B57}

\titleformat{\section}{\normalfont\Large\bfseries\color{matd48azul}}{\thesection}{1em}{}
\titleformat{\subsection}{\normalfont\large\bfseries\color{matd48azul}}{\thesubsection}{1em}{}

\providecommand{\ListaTitulo}{Lista de Exercícios}

\setlength{\headheight}{14pt}
\pagestyle{fancy}
\fancyhf{}
\lhead{\small MATD48 --- Planejamento de Experimentos A}
\rhead{\small \ListaTitulo}
\cfoot{\thepage}
\renewcommand{\headrulewidth}{0.4pt}

\lstset{
  basicstyle=\ttfamily\small,
  breaklines=true,
  frame=single,
  columns=fullflexible,
  backgroundcolor=\color{gray!8},
  commentstyle=\color{matd48verde},
  keywordstyle=\color{matd48azul}\bfseries,
  language=R,
  extendedchars=true,
  literate=%
    {á}{{\'a}}1 {à}{{\`a}}1 {â}{{\^a}}1 {ã}{{\~a}}1
    {é}{{\'e}}1 {ê}{{\^e}}1 {í}{{\'i}}1 {ó}{{\'o}}1
    {ô}{{\^o}}1 {õ}{{\~o}}1 {ú}{{\'u}}1 {ç}{{\c c}}1
    {Á}{{\'A}}1 {À}{{\`A}}1 {Â}{{\^A}}1 {Ã}{{\~A}}1
    {É}{{\'E}}1 {Ê}{{\^E}}1 {Í}{{\'I}}1 {Ó}{{\'O}}1
    {Ô}{{\^O}}1 {Õ}{{\~O}}1 {Ú}{{\'U}}1 {Ç}{{\c C}}1
}

\hypersetup{colorlinks=true, linkcolor=matd48azul, urlcolor=matd48azul, citecolor=matd48azul}

\newcommand{\cabecalho}[2]{%
  \begin{center}
    {\Large\bfseries\color{matd48azul} #1}\\[0.3em]
    {\large #2}\\[0.3em]
    \rule{\linewidth}{0.6pt}
  \end{center}
  \vspace{0.5em}
}

\newenvironment{questao}[1]{\vspace{1em}\noindent\textbf{Questão #1.}\hspace{0.4em}}{\par}
\newenvironment{solucao}{\vspace{0.3em}\noindent\textit{Solução.}\hspace{0.4em}\color{black!85}}{\par\vspace{0.5em}}


\begin{document}

\cabecalho{Gabarito --- Lista de Exercícios 12}{Fatoriais A×B, A×B×C e blocos completos (Capítulo 5 / Aula 12)}

\begin{questao}{1} \textbf{(Ciência de dados --- ANOVA de um fatorial A×B na mão)}
\end{questao}
\begin{solucao}
\begin{enumerate}[label=(\alph*)]
  \item Médias de célula: $\bar y_{LRU,peq}=12$, $\bar y_{LRU,gr}=22$, $\bar y_{LFU,peq}=18$,
    $\bar y_{LFU,gr}=36$. Marginais de $A$: $\bar y_{LRU\cdot}=17$, $\bar y_{LFU\cdot}=27$.
    Marginais de $B$: $\bar y_{\cdot peq}=15$, $\bar y_{\cdot gr}=29$. Média geral $\bar
    y_{\cdot\cdot}=22$.
  \item
  $$SQ_A = (2)(2)\big[(17-22)^2+(27-22)^2\big] = 4(25+25) = 200,$$
  $$SQ_B = (2)(2)\big[(15-22)^2+(29-22)^2\big] = 4(49+49) = 392,$$
  Resíduos de interação: $(12-17-15+22)=2$, $(22-17-29+22)=-2$, $(18-27-15+22)=-2$,
  $(36-27-29+22)=2$ (todos $\pm2$), logo
  $$SQ_{AB} = 2(2^2+2^2+2^2+2^2) = 2(16) = 32.$$
  Dentro de cada célula, os dois valores distam $\pm2$ da média da célula, então cada célula
  contribui $2^2+2^2=8$ para o erro: $SQ_E = 4(8) = 32$.

  \begin{center}
  \begin{tabular}{lccccc}
  \toprule
  Fonte & gl & SQ & QM & $F$ \\ \midrule
  $A$ & 1 & 200 & 200 & 25 \\
  $B$ & 1 & 392 & 392 & 49 \\
  $A{\times}B$ & 1 & 32 & 32 & 4 \\
  Erro & 4 & 32 & 8 & \\
  Total & 7 & 656 & & \\ \bottomrule
  \end{tabular}
  \end{center}
  \item Como $F_{1,4}=7{,}71$: $A$ ($F=25$) e $B$ ($F=49$) são significativos; $A{\times}B$
    ($F=4$) \textbf{não} é.
  \item O algoritmo de cache e o tamanho do payload afetam o tempo de resposta de forma
    aproximadamente \textbf{aditiva}: o efeito de trocar o algoritmo (LRU $\to$ LFU) é parecido
    em magnitude para os dois tamanhos de payload (10 ms para payload pequeno, $22-12=10$; 14 ms
    para payload grande, $36-22=14$ --- próximos, e a diferença não é grande o bastante para ser
    estatisticamente distinguível do ruído com $n=8$).
\end{enumerate}
\end{solucao}

\begin{questao}{2} \textbf{(Agricultura --- interpretando uma ANOVA real)}
\end{questao}
\begin{solucao}
\begin{enumerate}[label=(\alph*)]
  \item $F_{Irrigação}=5{,}28 > 4{,}49$: significativo. $F_{Silício}=9{,}38>3{,}24$:
    significativo. $F_{Irrigação:Silício}=2{,}63<3{,}24$: \textbf{não} significativo.
  \item Sugere que os três blocos de terreno eram, na prática, bastante \textbf{homogêneos} entre
    si: a variação real entre blocos era pequena, então "gastar" graus de liberdade do erro com o
    termo de bloco não trouxe ganho de precisão --- na verdade, piorou ligeiramente $QM_E$ porque
    reduziu os graus de liberdade do erro sem remover variação real.
  \item Sim. A decisão de blocar deve ser tomada a partir do \textbf{conhecimento do
    delineamento} (sabe-se, a priori, que parcelas de terrenos diferentes podem ter fertilidade
    diferente) e \emph{antes} de ver os dados --- não é uma escolha post-hoc guiada pelo
    resultado. Quando os blocos de fato capturam heterogeneidade real, o ganho é substancial
    (Capítulo 4); quando não, como aqui, o custo é pequeno, mas blocar continua sendo a decisão
    metodologicamente defensável dado o que se sabia no momento do planejamento.
\end{enumerate}
\end{solucao}

\begin{questao}{3} \textbf{(Dedução --- ortogonalidade no fatorial 2×2)}
\end{questao}
\begin{solucao}
\begin{enumerate}[label=(\alph*)]
  \item As quatro combinações de $(x_A,x_B)$, cada uma repetida $r$ vezes, são $(-1,-1)$,
    $(-1,+1)$, $(+1,-1)$, $(+1,+1)$. O produto $x_Ax_B$ vale, respectivamente, $+1,-1,-1,+1$.
    Somando sobre as $4r$ observações balanceadas:
    $$\mathbf{x}_A'\mathbf{x}_B = r(+1) + r(-1) + r(-1) + r(+1) = 0.$$
  \item Para $\mathbf{x}_A'\mathbf{x}_{AB}$, o produto $x_A \cdot (x_Ax_B) = x_A^2 x_B = x_B$
    (pois $x_A^2=1$). A soma de $x_B$ sobre as quatro combinações, cada uma $r$ vezes, é
    $r(-1)+r(+1)+r(-1)+r(+1)=0$. Por simetria, $\mathbf{x}_B'\mathbf{x}_{AB} =
    \sum x_B\cdot(x_Ax_B) = \sum x_A x_B^2 = \sum x_A = 0$ também.
  \item As três covariâncias cruzadas (fora da diagonal, excluindo o intercepto) são zero pelos
    itens (a)-(b); com a coluna de 1's do intercepto, $\mathbf{1}'\mathbf{x}_A =
    \mathbf{1}'\mathbf{x}_B = \mathbf{1}'\mathbf{x}_{AB}=0$ também (cada coluna soma zero, por
    construção balanceada). Logo $\mathbf{X}'\mathbf{X} = \mathrm{diag}(4r, 4r, 4r, 4r)$ ---
    diagonal.
  \item Com $\mathbf{X}'\mathbf{X}$ diagonal, $\hat\beta_{AB} =
    \mathbf{x}_{AB}'\mathbf{Y}/\mathbf{x}_{AB}'\mathbf{x}_{AB}$. O denominador é $4r$ (soma de
    $4r$ termos $(\pm1)^2=1$). O numerador, agrupando por célula ($x_{AB}=+1$ nas células
    $(1,1)$ e $(2,2)$; $x_{AB}=-1$ em $(1,2)$ e $(2,1)$):
    $$\mathbf{x}_{AB}'\mathbf{Y} = r\bar y_{11\cdot} - r\bar y_{12\cdot} - r\bar y_{21\cdot} +
    r\bar y_{22\cdot}.$$
    Logo
    $$\hat\beta_{AB} = \frac{r(\bar y_{11\cdot}-\bar y_{12\cdot}-\bar y_{21\cdot}+\bar
    y_{22\cdot})}{4r} = \frac{1}{4}\big(\bar y_{11\cdot}-\bar y_{12\cdot}-\bar y_{21\cdot}+\bar
    y_{22\cdot}\big),$$
    ou seja, a constante de proporcionalidade é $\boldsymbol{1/4}$ --- e, notavelmente, \textbf{não
    depende de $r$}: o número de repetições cancela entre numerador e denominador.
\end{enumerate}
\end{solucao}

\begin{questao}{4} \textbf{(Ciência de dados --- leitura causal de um teste A/B/n)}
\end{questao}
\begin{solucao}
\begin{enumerate}[label=(\alph*)]
  \item Médias marginais: azul $= (3{,}0+5{,}0)/2 = 4{,}0$; laranja $=(5{,}0+5{,}0)/2=5{,}0$.
    Efeito principal (laranja $-$ azul) $= 5{,}0-4{,}0 = \mathbf{1{,}0}$ ponto percentual.
  \item No celular: laranja $-$ azul $=5{,}0-3{,}0=2{,}0$. No desktop: laranja $-$ azul
    $=5{,}0-5{,}0=0{,}0$. Os dois valores \textbf{não} coincidem com o efeito principal
    marginal (1,0) --- um é o dobro, o outro é nulo.
  \item Há \textbf{heterogeneidade de efeito de tratamento} clara: o botão laranja só ajuda no
    celular; no desktop, a cor não faz diferença nenhuma. O efeito principal marginal (1,0 ponto)
    é uma média que \emph{não representa bem} nenhum dos dois subgrupos --- reportá-lo sozinho
    esconderia que a intervenção só funciona pela metade do público. A decisão de produto sugerida
    é \textbf{personalizar}: usar laranja para usuários de celular e manter qualquer cor (a
    diferença é nula) no desktop, em vez de escolher uma única cor "vencedora" para todos.
\end{enumerate}
\end{solucao}

\begin{questao}{5} \textbf{(Ciência de dados --- exercício computacional em R)}
\end{questao}
\begin{solucao}
\begin{enumerate}[label=(\alph*)]
  \item
  \begin{lstlisting}
X <- model.matrix(~ algoritmo * payload, data = dados)
identical(X[, 4], X[, 2] * X[, 3])
modelo <- aov(tempo ~ algoritmo * payload, data = dados)
summary(modelo)
  \end{lstlisting}
  \item Graus de liberdade: $A$ (algoritmo) tem 1 gl, $B$ (payload) tem 1 gl, $A{:}B$ tem 1 gl, e
    Resíduos têm $8 - 1 - 1 - 1 - 1 = 4$ gl (8 observações menos 1 para cada um dos 3 efeitos e 1
    para o intercepto) --- os mesmos graus de liberdade da tabela montada à mão na Questão 1(b),
    e os valores de $F$ reportados por \texttt{summary(modelo)} devem coincidir exatamente com
    $F_A=25$, $F_B=49$, $F_{AB}=4$.
\end{enumerate}
\end{solucao}

\end{document}
